\documentclass[10pt,a4paper]{article}
\usepackage[margin=1.6cm]{geometry}
\usepackage[T1]{fontenc}
\usepackage{lmodern}
\usepackage{microtype}
\usepackage{booktabs}
\usepackage{tabularx}
\usepackage{enumitem}
\usepackage{xcolor}
\usepackage[hidelinks]{hyperref}
\setlist{nosep,leftmargin=1.2em}
\setlength{\parskip}{3pt}
\setlength{\parindent}{0pt}
\newcommand{\code}[1]{\texttt{\small #1}}

\title{\vspace{-1.2cm}\textbf{Owl-Merge: Simple vs.\ Full Ontology Merge,\\ GraphRAG Evaluation and a Multi-hop SPARQL Question-Answering Engine}}
\author{Taher Merchant \quad\small\url{https://github.com/TaherMerchant25/Owl-Merge}}
\date{September 2026}

\begin{document}
\maketitle
\vspace{-0.9cm}

\section{Goal}
Biomedical knowledge is spread over many databases (Hetionet, DRKG, PrimeKG, OptimusKG, OBO ontologies, DBpedia, NCBI Gene), each naming the same gene or drug with a different IRI.
I received one pharmacological knowledge graph integrated two ways, as OWL files, and asked: \emph{does the merge strategy change what an AI system can correctly answer?}
The strategies follow Osman, Ben~Yahia \& Diallo, \emph{Ontology integration: approaches and challenging issues}, Information Fusion (2021), DOI \href{https://doi.org/10.1016/j.inffus.2021.01.007}{10.1016/j.inffus.2021.01.007}:
\textbf{Simple Merge} $O_3 = O_1 \cup O_2 \cup O_A$ keeps equivalent entities separate and links them with \code{owl:equivalentClass} axioms carrying a confidence score;
\textbf{Full Merge} $O_3 = O_1 \cup O_2$ fuses them into one entity that keeps one IRI and adds the others as labels. The paper argues both are semantically equivalent.

\section{What I built, step by step}
\begin{enumerate}
\item \textbf{Streaming OWL extractor} (\code{owl\_extract.py}). The files are 954\,MB and 659\,MB, too big for in-memory libraries such as rdflib. A line-oriented regex parser streams RDF/XML into entity, edge and alignment tables in about 50\,s.
\item \textbf{Structural analysis} with Python, \code{grep}/\code{awk}: entity and edge counts, equivalence links, class hierarchy, and a same-source ID check (a node holding two IDs from the same database, e.g.\ \code{Gene::675} and \code{Gene::5888}).
\item \textbf{GraphRAG evaluation with WeamRAG}, a hierarchical GraphRAG framework (GMM clustering, LLM community summaries, beam search). \code{build\_weamrag\_input.py} samples matched 3,000-node subgraphs from the same seed drugs and genes and converts them to WeamRAG's JSONL input; entity names carry their source (\code{HLA-A [hetionet:3105]}) so WeamRAG does not silently re-merge the Simple graph. Both indices were built locally (Llama~3.1~8B via Ollama, nomic-embed-text, Milvus Lite, SQLite).
\item \textbf{Question-answering test} (\code{compare\_merge\_rag.py}). The first version failed: embedding-only anchoring reached the target entities in 1 of 12 retrievals. I added entity-anchored seeding, a 10k-character context budget, reasoning-before-verdict prompts, two never-fused control pairs and a no-context control.
\item \textbf{Multi-hop SPARQL engine} (\code{kg\_engine.py}). Both OWL files are bulk-loaded into embedded on-disk \textbf{Oxigraph} stores (\code{pyoxigraph}); every graph access is a SPARQL~1.1 query.
\item \textbf{Web application} (\code{web\_app.py}, \code{web/index.html}). A dependency-free Python \code{http.server} backend with a validated JSON API and a plain HTML/CSS/JS page to ask a question against Simple, Full, or both side by side.
\end{enumerate}

\section{How the SPARQL engine answers a question}
\begin{enumerate}
\item \textbf{Retrieval}: question phrases are matched to \code{rdfs:label}s, longest n-gram first, with a single-word fallback (``tamoxifen'' $\rightarrow$ ``tamoxifen citrate'').
\item \textbf{Identity}: \code{owl:equivalentClass} links, plain or reified \code{owl:Axiom} with \code{pkg:confidenceScore}, are followed only above a threshold (default 0.9). This is how Simple Merge crosses databases.
\item \textbf{Reasoning}: two entities $\rightarrow$ bidirectional shortest-path search ($\le$3 relation hops); one entity $\rightarrow$ hop-by-hop expansion where type words (``drugs \ldots diseases'') fix the target type per hop and relations are ranked by embedding similarity to the question (e.g.\ \emph{indication} 0.58 followed, \emph{contraindication} 0.50 skipped). The trace is shown.
\item \textbf{Answer}: Llama~3.1~8B answers only from numbered evidence paths and cites them (\code{[P1]}).
\item \textbf{Safety checks}: a \emph{WARNING} for nodes fusing two same-source IDs and a \emph{NOTE} when two named entities resolve to one node by label alone.
\end{enumerate}

\section{Results}
\begin{center}
\small
\begin{tabular}{@{}lrr@{}}
\toprule
Measure & Simple Merge & Full Merge \\
\midrule
Declared entities & 268,654 & 97,581 \\
Edges & 3,375,438 & 2,531,908 \\
\code{equivalentClass} links (scored / unscored) & 246,934 (208,688 / 38,246) & 0 \\
Classes with multiple parents & 36.5\% & 35.2\% \\
Entities holding $\ge$2 IDs from the same source & \textbf{0} & \textbf{96} \\
Relations in a matched 3,000-node sample & 38,244 & 114,164 \\
Triples in Oxigraph store (size) & 5,343,202 (421\,MB) & 3,065,773 (238\,MB) \\
\bottomrule
\end{tabular}
\end{center}
\begin{itemize}
\item \textbf{Full Merge is 63.7\% smaller but fuses different entities.} All 7 hand-checked same-source fusions are real errors: BRCA2+RAD51, atracurium+cisatracurium, estrone+estropipate, LGALS2+LGALS3, HLA-A+HLA-E, SYT11+SYT12, CERS1+GDF1.
\item \textbf{Errors enter through identity chains.} LGALS2 was linked to Wikidata's LGALS3 at 0.871; HLA-A reaches HLA-E in three links (weakest 0.701). The engine found BRCA2 $\equiv$ \code{dbpedia:RAD51} (unscored) $\equiv$ RAD51 (0.72).
\item \textbf{GraphRAG passes the error to the LLM.} With no context the model judged all 7 pairs \emph{different}. With Full Merge context it answered \emph{same} for 4 of 5 fused pairs, quoting the fused node as evidence; with Simple Merge context it said \emph{same} once. Small test (7 pairs, one 8B model), so suggestive, not conclusive.
\item \textbf{The engine shows the difference directly.} ``Are BRCA2 and RAD51 the same gene?'' returns \code{BRCA2 --protein protein--> RAD51} (they interact) on Simple Merge, but only a fused node plus a WARNING on Full Merge. ``How is BRCA2 related to tamoxifen?'' crosses three databases via 0.95 identity links on Simple Merge. Queries take seconds without the LLM.
\end{itemize}

\section{Technology used}
\begin{tabularx}{\linewidth}{@{}lX@{}}
\toprule
Data \& standards & OWL, RDF/XML, SPARQL 1.1, reified \code{owl:Axiom} alignments \\
Triple store & Oxigraph (\code{pyoxigraph} 0.5.11), embedded, on-disk \\
GraphRAG & WeamRAG, Milvus Lite, SQLite, UMAP/GMM clustering \\
Models (local) & Ollama: Llama 3.1 8B (answers), nomic-embed-text (relation ranking); NVIDIA TITAN RTX \\
Backend \& frontend & Python 3 standard library (\code{http.server}), HTML/CSS/JavaScript, no build step \\
Tooling & Git/GitHub, Playwright (UI testing), VS Code Remote-SSH \\
\bottomrule
\end{tabularx}

\section{How to run}
\begin{verbatim}
python3 -m venv .venv && .venv/bin/pip install pyoxigraph==0.5.11
ollama pull llama3.1:8b && ollama pull nomic-embed-text
.venv/bin/python kg_engine.py load simple && .venv/bin/python kg_engine.py load full
.venv/bin/python kg_engine.py ask simple "How is BRCA2 related to tamoxifen?"
.venv/bin/python web_app.py        # open http://127.0.0.1:8780
\end{verbatim}

\section{Limitations and next steps}
Only the first two named entities in a question are connected; name matching is exact or single-word, so misspellings (e.g.\ ``metformonin'') do not match; target types come from IRI patterns; the QA test is small. The two OWL files were not built from exactly identical inputs, so some differences reflect inputs rather than strategy.
Next: review the remaining 89 same-source fusions, add fuzzy entity matching, and evaluate on a larger question set with a stronger model.
\end{document}
